% Visual reference for every core data structure used in this volume.

\clearpage
\phantomsection
\addcontentsline{toc}{section}{Visual Data Structure Guide}
\section*{Visual Data Structure Guide}
\markboth{Visual Data Structure Guide}{Foundations}

Algorithms are easier to follow when saved values have a clear shape. Orange
shows what we are working on now, teal shows what is finished or kept, and
navy shows saved information. The different borders also remain clear when
the page is printed in grayscale.

% A guide entry should never leave its heading and explanation at the bottom
% while moving the diagram to the next page.
\newcommand{\visualguidesection}[1]{\Needspace{22\baselineskip}\subsection{#1}}

\begin{visualbox}{Visual legend used throughout the book}
\centering
\begin{tikzpicture}[font=\small,>={Stealth[length=2.2mm,width=1.55mm]},node distance=13mm]
  \node[dsnode] (stored) {saved};
  \node[dsnode,active,right=of stored] (current) {working now};
  \node[dsnode,visited,right=of current] (done) {finished / kept};
  \draw[flowarrow] (stored)--(current);
  \draw[flowarrow,teal] (current)--(done);
\end{tikzpicture}
\end{visualbox}

\visualguidesection{Array and \texttt{vector}}
An array is a row of values. Each box has a number called an index. A C++
\texttt{vector} can grow, but we still use indexes in the same way.

\begin{visualbox}{An array is a row of numbered boxes}
\centering
\begin{tikzpicture}[font=\small,>={Stealth[length=2.2mm,width=1.55mm]}]
  \foreach \x/\v in {0/4,1/7,2/1,3/9,4/3,5/6} {
    \node[arraycell] (a\x) at (1.15*\x,0) {\v};
    \node[below=1mm of a\x,font=\scriptsize,text=navy] {\x};
  }
  \node[above=8mm of a0,text=teal,font=\bfseries] (l) {left};
  \node[above=8mm of a5,text=orange,font=\bfseries] (r) {right};
  \draw[->,very thick,teal] (l)--(a0.north);
  \draw[->,very thick,orange] (r)--(a5.north);
  \draw[decorate,decoration={brace,amplitude=5pt},navy]
    ([yshift=-8mm]a0.south west)--([yshift=-8mm]a5.south east)
    node[midway,below=6mm]{use index \texttt{i} to open box \texttt{values[i]}};
\end{tikzpicture}
\end{visualbox}

Arrays appear in two-pointer scans, binary search, bit results, prices, house
values, traversal orders, and one-dimensional dynamic programming.

\visualguidesection{String and reading position}
A string is a row of characters. We keep one position to show where reading
starts and use the length to know where it stops.

\begin{visualbox}{The first number tells us how many letters to read}
\centering
\begin{tikzpicture}[font=\small,>={Stealth[length=2.2mm,width=1.55mm]}]
  \foreach \x/\v in {0/4,1/\#,2/c,3/o,4/d,5/e,6/3,7/\#,8/c,9/+,10/+} {
    \node[arraycell,minimum width=8.5mm] (s\x) at (0.88*\x,0) {\v};
    \node[below=1mm of s\x,font=\scriptsize,text=navy] {\x};
  }
  \node[dsnode,active,above=9mm of s0] (cursor) {start here};
  \draw[flowarrow,orange] (cursor)--(s0.north);
  \node[dsnode,visited,above=9mm of s3] (payload) {read the next 4};
  \draw[flowarrow,teal] (payload)--(s3.north);
  \draw[decorate,decoration={brace,amplitude=5pt},teal]
    ([yshift=-8mm]s2.south west)--([yshift=-8mm]s5.south east)
    node[midway,below=6mm]{the word \texttt{code}};
\end{tikzpicture}

\smallskip
The first number is 4, so we read the next four letters.
\end{visualbox}

\visualguidesection{Remember the best earlier value}
For stock profit, we only need the cheapest price seen before today.

\begin{visualbox}{For each sell day, use the cheapest earlier price}
\centering
\begin{tikzpicture}[font=\small,>={Stealth[length=2.2mm,width=1.55mm]}]
  \foreach \x/\v in {0/7,1/1,2/5,3/3,4/6,5/4} {
    \node[arraycell,minimum width=11mm] (p\x) at (1.18*\x,0) {\v};
    \node[below=1mm of p\x,font=\scriptsize] {day \x};
  }
  \node[dsnode,visited,above=11mm of p1] (buy) {cheapest earlier price\\buy at 1};
  \node[dsnode,active,above=11mm of p4] (sell) {sell now at 6};
  \draw[flowarrow,teal] (buy)--(p1.north);
  \draw[flowarrow,orange] (sell)--(p4.north);
  \draw[->,very thick,navy] (p1.north east) to[bend left=24]
    node[edgelabel,above]{profit \(6-1=5\)} (p4.north west);
\end{tikzpicture}
\end{visualbox}

\visualguidesection{Hash set and hash map}
A hash set answers ``Have I seen this key?'' A hash map saves extra information
beside the key, such as an index.

\begin{visualbox}{A key is changed into a small bucket number}
\centering
\begin{tikzpicture}[font=\small,>={Stealth[length=2.2mm,width=1.55mm]},node distance=5mm]
  \node[dsnode,fill=softorange] (key) {key 7};
  \node[dsnode,right=14mm of key] (hash) {hash\\rule};
  \draw[flowarrow] (key)--(hash);
  \foreach \y/\label in {0/0,1/1,2/2,3/3} {
    \node[draw=navy,minimum width=18mm,minimum height=7mm,
      right=18mm of hash,yshift={(1.5-\y)*8mm}] (b\y) {bucket \label};
  }
  \draw[flowarrow,orange] (hash.east)--(b3.west);
  \node[dsnode,right=14mm of b3,fill=softteal] (set) {set:\\found 7};
  \node[dsnode,right=8mm of set,fill=softblue] (map) {map:\\saved index 4};
  \draw[flowarrow,teal] (b3)--(set);
  \draw[flowarrow] (set)--(map);
\end{tikzpicture}

\smallskip
A hash table normally finds a saved key quickly. Spreading keys among the
buckets keeps it fast.
\end{visualbox}

\visualguidesection{Stack and function calls}
A stack handles the newest item first. Function calls use a stack too: the
newest call finishes before the older call below it.

\begin{visualbox}{Function calls stack on top of one another}
\centering
\begin{tikzpicture}[font=\small,>={Stealth[length=2.2mm,width=1.55mm]}]
  \node[draw=navy,fill=softblue,minimum width=42mm,minimum height=8mm] (f0) at (0,0)
    {\texttt{dfs(node 3)}};
  \node[draw=navy,fill=softblue,minimum width=42mm,minimum height=8mm,
    above=0mm of f0] (f1) {\texttt{dfs(node 20)}};
  \node[draw=orange,fill=softorange,minimum width=42mm,minimum height=8mm,
    above=0mm of f1] (f2) {\texttt{dfs(node 15)}};
  \node[left=14mm of f2,text=orange,font=\bfseries,align=right] (top) {next to\\finish};
  \draw[flowarrow,orange] (top)--(f2);
  \draw[->,thick,teal] ([xshift=12mm]f2.east) arc[start angle=90,end angle=-90,radius=9mm]
    node[edgelabel,right=7mm,midway,align=left]{finish\\and remove};
\end{tikzpicture}
\end{visualbox}

The stack grows by one box for each unfinished function call.

\visualguidesection{A problem can split into choices}
This picture shows every choice a function may try. Each choice creates a
smaller version of the same problem.

\begin{visualbox}{One problem splits into smaller choices}
\centering
\begin{tikzpicture}[font=\small,>={Stealth[length=2.2mm,width=1.55mm]}]
  \node[dsnode,active] (r) {solve this case};
  \node[dsnode,below left=13mm and 22mm of r] (a) {try choice A};
  \node[dsnode,below right=13mm and 22mm of r] (b) {try choice B};
  \node[dsnode,visited,below left=12mm and 7mm of a] (a1) {small case};
  \node[dsnode,below right=12mm and 7mm of a] (a2) {smaller case};
  \node[dsnode,below=12mm of b] (b1) {smaller case};
  \draw[flowarrow] (r)--(a); \draw[flowarrow] (r)--(b);
  \draw[flowarrow] (a)--(a1); \draw[flowarrow] (a)--(a2);
  \draw[flowarrow] (b)--(b1);
  \draw[flowarrow,teal] (a1.west)
    to[out=180,in=180,looseness=1.35]
    node[edgelabel,left,font=\footnotesize]{back} (a.west);
\end{tikzpicture}
\end{visualbox}

\visualguidesection{Queue}
A queue removes the oldest item first. We remove from the front and add new
items at the back.

\begin{visualbox}{Queue: remove at the front, add at the back}
\centering
\begin{tikzpicture}[font=\small,>={Stealth[length=2.2mm,width=1.55mm]}]
  \foreach \x/\v in {0/9,1/20,2/15,3/7} {
    \node[arraycell,minimum width=13mm] (q\x) at (1.35*\x,0) {\v};
  }
  \node[left=10mm of q0,text=teal,font=\bfseries] (front) {front};
  \node[right=10mm of q3,text=orange,font=\bfseries] (back) {back};
  \draw[flowarrow,teal] (front)--(q0);
  \draw[flowarrow,orange] (back)--(q3);
  \node[above=9mm of q0,text=teal,font=\bfseries] (pop) {remove next};
  \node[above=9mm of q3,text=orange,font=\bfseries] (push) {add new};
  \draw[flowarrow,teal] (q0.north)--(pop.south);
  \draw[flowarrow,orange] (push.south)--(q3.north);
\end{tikzpicture}
\end{visualbox}

\visualguidesection{Singly linked list}
Each node stores a value and an arrow to the next node. To move through the
list, follow the arrows.

\begin{visualbox}{Each node points to the next node}
\centering
\begin{tikzpicture}[font=\small,>={Stealth[length=2.2mm,width=1.55mm]}]
  \foreach \x/\v in {0/3,1/2,2/0,3/-4} {
    \node[draw=navy,fill=softblue,minimum width=9mm,minimum height=9mm] (v\x)
      at (2.45*\x,0.7) {\v};
    \node[draw=navy,fill=white,minimum width=11mm,minimum height=9mm,
      right=0mm of v\x,font=\scriptsize] (p\x) {next};
  }
  \draw[flowarrow] (p0.east)--(v1.west);
  \draw[flowarrow] (p1.east)--(v2.west);
  \draw[flowarrow] (p2.east)--(v3.west);
  \draw[->,very thick,orange] (p3.south)
    to[out=-90,in=-90,looseness=1.35]
    node[edgelabel,midway,below=9mm,font=\small]{last node points back to 2} (v1.south);
  \node[above=7mm of v0,text=teal,font=\bfseries] (head) {head};
  \draw[flowarrow,teal] (head)--(v0.north);
\end{tikzpicture}
\end{visualbox}

\visualguidesection{Binary tree}
A binary-tree node can point to a left child and a right child. Every node
below one node forms that node's branch.

\begin{visualbox}{A node can have a left and a right child}
\centering
\begin{tikzpicture}[font=\small,>={Stealth[length=2.2mm,width=1.55mm]}]
  \node[trienode,active] (r) {3};
  \node[trienode,below left=13mm and 22mm of r] (l) {9};
  \node[trienode,below right=13mm and 22mm of r] (q) {20};
  \node[trienode,visited,below left=13mm and 10mm of q] (a) {15};
  \node[trienode,visited,below right=13mm and 10mm of q] (b) {7};
  \draw[flowarrow] (r)--node[edgelabel,above left,font=\scriptsize]{left} (l);
  \draw[flowarrow] (r)--node[edgelabel,above right,font=\scriptsize]{right} (q);
  \draw[flowarrow] (q)--(a); \draw[flowarrow] (q)--(b);
  \node[draw=teal,dashed,fit=(q)(a)(b),inner sep=4mm,
    label=right:{\scriptsize this whole branch starts at 20}] {};
\end{tikzpicture}
\end{visualbox}

\visualguidesection{Trie: words that share a beginning}
A Trie stores one letter at a time. Words with the same beginning follow the
same path until their letters become different.

\begin{visualbox}{The words \texttt{app} and \texttt{apple} share one path}
\centering
\begin{tikzpicture}[font=\small,>={Stealth[length=2.2mm,width=1.55mm]},node distance=11mm]
  \node[trienode] (r) {root};
  \node[trienode,right=of r] (a) {a};
  \node[trienode,right=of a] (p1) {p};
  \node[trienode,visited,right=of p1] (p2) {p};
  \node[trienode,right=of p2] (l) {l};
  \node[trienode,visited,right=of l] (e) {e};
  \draw[flowarrow] (r)--(a)--(p1)--(p2)--(l)--(e);
  \node[below=5mm of p2,text=teal]{\texttt{app} ends here};
  \node[below=5mm of e,text=teal]{\texttt{apple} ends here};
\end{tikzpicture}
\end{visualbox}

\visualguidesection{Graph and adjacency list}
A graph has numbered points joined by lines. For each point, a neighbor list
records the points directly connected to it.

\begin{visualbox}{The same graph shown as neighbor lists}
\centering
\begin{tikzpicture}[font=\small,>={Stealth[length=2.2mm,width=1.55mm]}]
  \node[trienode] (g1) at (0,1) {1};
  \node[trienode] (g2) at (1.6,1) {2};
  \node[trienode] (g3) at (1.6,-0.5) {3};
  \node[trienode] (g4) at (0,-0.5) {4};
  \draw[-,thick,navy] (g1)--(g2)--(g3)--(g4)--(g1);
  \node[dsnode,align=left,right=25mm of g2] (adj) {
    \textbf{neighbors of each point}\\
    1: [2,4]\\2: [1,3]\\3: [2,4]\\4: [1,3]
  };
  \draw[flowarrow,teal] (g2.east)--(adj.west);
\end{tikzpicture}
\end{visualbox}

\visualguidesection{Heap: quickly see the largest or smallest value}
A max heap keeps the largest value on top. A min heap keeps the smallest value
on top, so that value can be taken next.

\begin{visualbox}{The top shows the next value we can take}
\centering
\begin{tikzpicture}[font=\small,>={Stealth[length=2.2mm,width=1.55mm]}]
  \node[trienode,active] (ma) at (0,1.4) {5};
  \node[trienode,below left=10mm and 9mm of ma] (mb) {2};
  \node[trienode,below right=10mm and 9mm of ma] (mc) {3};
  \draw[-,thick,navy] (ma)--(mb); \draw[-,thick,navy] (ma)--(mc);
  \node[above=4mm of ma,text=orange,font=\bfseries]{largest on top};

  \node[trienode,visited] (na) at (5.2,1.4) {6};
  \node[trienode,below left=10mm and 9mm of na] (nb) {8};
  \node[trienode,below right=10mm and 9mm of na] (nc) {9};
  \draw[-,thick,navy] (na)--(nb); \draw[-,thick,navy] (na)--(nc);
  \node[above=4mm of na,text=teal,font=\bfseries]{smallest on top};
  \draw[<->,very thick,navy] (ma.east)--node[edgelabel,above]{middle values meet here} (na.west);
\end{tikzpicture}
\end{visualbox}

\visualguidesection{A table that reuses earlier answers}
A DP array saves answers already found. The colored boxes show which earlier
answers are used for the next box.

\begin{visualbox}{Choose the better of two earlier answers}
\centering
\begin{tikzpicture}[font=\small,>={Stealth[length=2.2mm,width=1.55mm]}]
  \foreach \x/\v in {0/0,1/2,2/7,3/11,4/11,5/12} {
    \node[arraycell] (d\x) at (1.15*\x,0) {\v};
    \node[below=1mm of d\x,font=\scriptsize] {\(dp[\x]\)};
  }
  \node[draw=teal,line width=1pt,fit=(d3),inner sep=0.6pt] {};
  \node[draw=orange,line width=1pt,fit=(d4),inner sep=0.6pt] {};
  \node[draw=orange,very thick,fit=(d5),inner sep=1mm] {};
  \node[dsnode,fill=softblue,right=16mm of d5,align=left] (rule) {
    \textcolor{orange}{option 1: }keep 11\\
    \textcolor{teal}{option 2: }\(11+1=12\)\\
    \textbf{choose: }12};
  \node[below=8mm of d5,text=orange,font=\bfseries] {answer for this box};
\end{tikzpicture}
\end{visualbox}

First decide what each box means. Then fill the boxes only after the earlier
answers they need are ready.

\visualguidesection{Binary-search interval}
Binary search keeps the part that may still contain the answer and throws away
the half that cannot contain it.

\begin{visualbox}{Keep only the half that may contain the answer}
\centering
\begin{tikzpicture}[font=\small,>={Stealth[length=2.2mm,width=1.55mm]}]
  \foreach \x/\v in {0/4,1/5,2/6,3/7,4/0,5/1,6/2} {
    \node[arraycell] (b\x) at (0.95*\x,0) {\v};
  }
  \node[above=8mm of b0,text=teal,font=\bfseries] (left) {left};
  \node[above=8mm of b3,text=orange,font=\bfseries] (mid) {mid};
  \node[above=8mm of b6,text=navy,font=\bfseries] (right) {right};
  \draw[flowarrow,teal] (left)--(b0.north);
  \draw[flowarrow,orange] (mid)--(b3.north);
  \draw[flowarrow] (right)--(b6.north);
  \node[draw=orange,very thick,dashed,fit=(b4)(b5)(b6),inner sep=2mm,
    label=below:{keep this half}] {};
\end{tikzpicture}
\end{visualbox}

\visualguidesection{Binary digits and right shift}
A number can be written with only 0s and 1s. Shifting right removes the last
binary digit and leaves a smaller number.

\begin{visualbox}{Shifting right drops the last binary digit}
\centering
\begin{tikzpicture}[font=\small,>={Stealth[length=2.2mm,width=1.55mm]}]
  \node[arraycell] (x3) at (0,0) {1};
  \node[arraycell,right=0mm of x3] (x2) {1};
  \node[arraycell,right=0mm of x2] (x1) {0};
  \node[arraycell,active,right=0mm of x1] (x0) {1};
  \node[left=10mm of x3,font=\bfseries,text=navy] {13};
  \node[right=9mm of x0,text=orange] {last bit is 1};

  \node[arraycell,below=15mm of x3] (y2) {1};
  \node[arraycell,right=0mm of y2] (y1) {1};
  \node[arraycell,right=0mm of y1] (y0) {0};
  \node[left=10mm of y2,font=\bfseries,text=teal] {6};
  \draw[flowarrow,teal] (x2.south)--node[edgelabel,right]{shift 13 right} (y1.north);
  \node[right=9mm of y0,text=navy] {13 becomes 6};
\end{tikzpicture}
\end{visualbox}

\clearpage